\documentclass[10pt]{article}
\usepackage[utf8]{inputenc}
\usepackage[T1]{fontenc}
\usepackage{amsmath}
\usepackage{amsfonts}
\usepackage{amssymb}
\usepackage[version=4]{mhchem}
\usepackage{stmaryrd}
\usepackage{bbold}
\usepackage{caption}
\usepackage{multirow}
\usepackage{hyperref}
\hypersetup{colorlinks=true, linkcolor=blue, filecolor=magenta, urlcolor=cyan,}
\urlstyle{same}

\title{A new almost perfect nonlinear function which is not quadratic }

\author{Yves Edel* Alexander Pott ${ }^{\dagger}$}
\date{}


%New command to display footnote whose markers will always be hidden
\let\svthefootnote\thefootnote
\newcommand\blfootnotetext[1]{%
  \let\thefootnote\relax\footnote{#1}%
  \addtocounter{footnote}{-1}%
  \let\thefootnote\svthefootnote%
}

%Overriding the \footnotetext command to hide the marker if its value is `0`
\let\svfootnotetext\footnotetext
\renewcommand\footnotetext[2][?]{%
  \if\relax#1\relax%
    \ifnum\value{footnote}=0\blfootnotetext{#2}\else\svfootnotetext{#2}\fi%
  \else%
    \if?#1\ifnum\value{footnote}=0\blfootnotetext{#2}\else\svfootnotetext{#2}\fi%
    \else\svfootnotetext[#1]{#2}\fi%
  \fi
}

\DeclareUnicodeCharacter{2218}{\ifmmode\circ\else{$\circ$}\fi}

\begin{document}
\maketitle
\captionsetup{singlelinecheck=false}


\begin{abstract}
Following an example in [11], we show how to change one coordinate function of an almost perfect nonlinear (APN) function in order to obtain new examples. It turns out that this is a very powerful method to construct new APN functions. In particular, we show that our approach can be used to construct a "non-quadratic" APN function. This new example is in remarkable contrast to all recently constructed functions which have all been quadratic. An equivalent function has been found independently by Brinkmann and Leander [7]. However, they claimed that their function is CCZ equivalent to a quadratic one. In this paper we give several reasons why this new function is not equivalent to a quadratic one.
\end{abstract}

\section*{1 Preliminaries}
We consider functions $F: \mathbb{F}_{2}{ }^{n} \rightarrow \mathbb{F}_{2}{ }^{n}$ with "good" differential and linear properties. Motivated by applications in cryptography, a lot of research has been done to construct functions which are "as nonlinear as possible". We discuss two possibilities to define nonlinearity: The first one uses differential properties, the second measures the "distance" to linear functions.

Let us begin with the differential properties. Given $F: \mathbb{F}_{2}{ }^{n} \rightarrow \mathbb{F}_{2}{ }^{n}$, we define

$$
\Delta_{F}(a, b):=|\{x: F(x+a)-F(x)=b\}| .
$$

We have $\Delta_{F}(0,0)=2^{n}$, and $\Delta_{F}(0, b)=0$ if $b \neq 0$.\\
Since we are working in fields of characteristic 2, we may replace the " - " by "+" and write $F(x+a)+F(x)$ instead of $F(x+a)-F(x)$. We say that $F$ is almost perfect nonlinear (APN) if $\Delta_{F}(a, b) \in\{0,2\}$ for all $a, b \in \mathbb{F}_{2}^{n}, a \neq 0$. Note that $\Delta_{F}(a, b) \in\left\{0,2^{n}\right\}$ if $F$ is linear, hence the condition $\Delta_{F}(a, b) \in\{0,2\}$ identifies functions which are quite different from linear mappings. In characteristic 2, it is impossible that $\Delta_{F}(a, b)=1$ : If $x$ is a solution of $F(x+a)-F(x)=b$, then $x+a$ is another solution, so that all values of $\Delta_{F}$ must be even. In the case of odd characteristic, functions $F: \mathbb{F}_{q}{ }^{n} \rightarrow \mathbb{F}_{q}{ }^{n}$ with $\Delta_{F}(a, b)=1$ for all $a \neq 0$ do exist; they are called perfect nonlinear or planar. In the last few years, many new APN functions have been constructed. The first example of a non-power mapping has been described in [27]. Infinite series are contained in $[5,10,11,13,16,17]$. Also some new planar functions have been found, see $[15,23,37]$.

There may be a possibility for a unified treatment of (some of) these constructions in the even and odd characteristic case. In particular, we suggest to look more carefully at the underlying design of an APN function, similar to the designs corresponding to planar functions, which are projective planes, see [30].

In this paper, we consider APN functions only. In Section 2, we discuss the problem how to distinguish APN functions up to equivalence. We discuss the concept of affine, extended affine and CCZ equivalence. We also introduce a certain design associated with an APN function, which

\footnotetext{*Department of Pure Mathematics and Computer Algebra, Ghent University, Krijgslaan 281, S22, B-9000 Ghent, Belgium. The research of the author takes place within the project "Linear codes and cryptography" of the Fund for Scientific Research Flanders (FWO-Vlaanderen) (Project nr. G.0317.06), and is supported by the Interuniversitary Attraction Poles Programme - Belgian State - Belgian Science Policy: project P6/26-Bcrypt.\\
${ }^{\dagger}$ Department of Mathematics, Otto-von-Guericke-University Magdeburg, D-39016 Magdeburg, Germany
}
can be used to distinguish inequivalent functions. In Theorem 7, we obtain a necessary condition for an APN function to be crooked (see Definition 4).

In Section 3, we describe an idea due to Dillon for obtaining new APN functions from known ones by changing a component function. We apply this switching construction to the known examples of APN functions in small dimensions ( $n \leq 9$ ). The results are tabulated at the end of the paper. In particular, we construct a new non-quadratic APN function (Theorem 11). We compute several parameters (for instance automorphism groups) which are invariant under CCZ equivalence. Section 4 contains some comments and problems concerning these invariants.

Another approach to nonlinearity is to measure the distance between linear functions $h$ : $\mathbb{F}_{2}{ }^{n} \rightarrow \mathbb{F}_{2}$ and the component functions $F_{g}: \mathbb{F}_{2}{ }^{n} \rightarrow \mathbb{F}_{2}$ : Component functions are defined via $F_{g}(x):=g(F(x))$ where $g$ is a nonzero linear function $\mathbb{F}_{2}{ }^{n} \rightarrow \mathbb{F}_{2}$. We denote the set of all linear functions $f: \mathbb{F}_{2}^{n} \rightarrow \mathbb{F}_{2}$ by $\widehat{\mathbb{F}_{2}^{n}}$. The Hamming distance $d_{H}(f, g)$ between two Boolean functions $f, g: \mathbb{F}_{2}{ }^{n} \rightarrow \mathbb{F}_{2}$ is simply the number of $x$ such that $f(x) \neq g(x)$. The Hamming weight $w_{H}(f)$ of a Boolean function $f$ is the number of $x$ with $f(x)=1$.

We say that a function $F$ is highly nonlinear if


\begin{equation*}
\min _{f, g \in \widehat{\mathbb{F}_{2}^{n}}, g \neq 0}\left(d_{H}\left(f, F_{g}\right), d_{H}\left(f+1, F_{g}\right)\right) \tag{1}
\end{equation*}


is large, i.e. the component functions $g \circ F=F_{g}$ of $F$ are as different as possible from all affine linear functions $f$ and $f+1$, where $f \in \widehat{\mathbb{F}_{2}^{n}}$ (here " ∘" denotes the composition of functions).

Instead of investigating $d_{H}\left(f, F_{g}\right)$ and $d_{H}\left(f+1, F_{g}\right)$, we may equivalently consider the Walsh coefficients

$$
\mathcal{W}_{F}(f, g)=\sum_{x \in \mathbb{F}_{2}^{n}}(-1)^{(g \circ F)(x)+f(x)}
$$

(sometimes we omit the subscript $F$ ). We have

$$
2^{n}-2 d_{H}\left(f, F_{g}\right)=\mathcal{W}_{F}(f, g) \quad \text { and } \quad 2^{n}-2 d_{H}\left(f+1, F_{g}\right)=-\mathcal{W}_{F}(f, g)
$$

This shows that the distances come in pairs $d_{1}$ and $d_{2}$ with $d_{1}+d_{2}=2^{n}$. Instead of maximizing the minimum of the $d_{H}\left(f, F_{g}\right), d_{H}\left(f+1, F_{g}\right)$ with $g \neq 0$, we may equivalently minimize the maximum of $|\mathcal{W}(f, g)|, g \neq 0$.

The Walsh coefficients are basically the weights of the following code of length $2^{n}$ : Let $F$ : $\mathbb{F}_{2}{ }^{n} \rightarrow \mathbb{F}_{2}{ }^{n}$ be any function. Define a matrix $C_{F} \in \mathbb{F}_{2}^{\left(2 n, 2^{n}\right)}$ as follows: The columns are the vectors $\binom{x}{F(x)}, x \in \mathbb{F}_{2}^{n}$. Then the rows of the matrix

$$
C_{F}=\left(\begin{array}{ccc}
\cdots & x & \cdots \\
\cdots & F(x) & \cdots
\end{array}\right)_{x \in \mathbb{F}_{2}^{n}}
$$

generate a binary code $\mathcal{C}_{F}$ whose codewords are the vectors

$$
v(f, g)=(f(x)+(g \circ F)(x))_{x \in \mathbb{F}_{2}^{n}}
$$

where $f$ and $g$ are linear functions $\mathbb{F}_{2}^{n} \rightarrow \mathbb{F}_{2}$. Note that the $2^{n}$ linear functions $\mathbb{F}_{2}^{n} \rightarrow \mathbb{F}_{2}$ can be described by the inner products $\langle c, x\rangle, c \in \mathbb{F}_{2}{ }^{n}$. Hence, the vectors $v(f, g)$ are indeed all the linear combinations of the rows of $C_{F}$.

It is easy to see that the Hamming weight $w_{H}(v(f, g))$ of the codeword $v(f, g)$ is related to the Walsh coefficient $\mathcal{W}(f, g)$ as follows:

$$
2^{n}-2 w_{H}(v(f, g))=\mathcal{W}(f, g)
$$

If the code $\mathcal{C}_{F}$ does not contain the vector $(1, \ldots, 1)$, we may add this vector as a row to $C_{F}$. The vector space generated by the rows of this extended matrix is called the extended code $\mathcal{C}_{F}^{\text {ext }}$\\
associated to the function $F$. This construction means that we add the vectors $w:=v(f+1, g)$ to the code $\mathcal{C}_{F}$. If $u:=v(f, g)$, we have $w_{H}(u)+w_{H}(w)=2^{n}$ and therefore

$$
2^{n}-2 w_{H}(u)=-\left(2^{n}-2 w_{H}(w)\right)
$$

which gives rise to the Walsh coefficients $\pm \mathcal{W}(f, g)$. We note that the vector ( $1, \ldots 1$ ) is not contained in $C_{F}$ if $F$ is APN, see [8, 19], for instance.

The multiset of values $\mathcal{W}_{F}(f, g)$ for all linear functions $f, g$ is called the Walsh spectrum of $F$.

The Walsh coefficients $\mathcal{W}_{F}(f, 0)$ are

$$
\mathcal{W}_{F}(f, 0)= \begin{cases}2^{n} & \text { if } f(x)=0 \text { for all } x \\ 0 & \text { otherwise }\end{cases}
$$

Therefore, the Walsh coefficients $\mathcal{W}_{F}(f, 0)$ do not depend on $F$. For this the reason, the Walsh coefficients with $g=0$ are sometimes excluded from the Walsh spectrum.

Usually, the Walsh spectrum is defined in terms of the trace function of a finite field. This "finite field definition" is completely equivalent to ours. We have used the vector space definition in order to emphasize that the Walsh spectrum (or the Walsh transformation) is just a property of the additive group of $\mathbb{F}_{2}^{n}$. If we identify $\mathbb{F}_{2}^{n}$ with the additive group of the finite field $\mathbb{F}_{2^{n}}$, then the linear mappings $f: \mathbb{F}_{2}{ }^{n} \rightarrow \mathbb{F}_{2}$ are just the mappings $f_{\alpha}$ defined by $x \mapsto \operatorname{tr}(\alpha x)$, where tr is the usual trace function $\operatorname{tr}(x):=\sum_{i=0}^{n-1} x^{2^{i}}$. We have $f_{\alpha} \neq f_{\beta}$ for $\alpha \neq \beta$, and we define

$$
\mathcal{W}_{F}\left(f_{\alpha}, f_{\beta}\right)=: \mathcal{W}_{F}(\alpha, \beta)
$$

We have $\mathcal{W}_{F}(\alpha, 0)=0$ if $\alpha \neq 0$ and $\mathcal{W}_{F}(0,0)=2^{n}$.\\
It is well known that there are $\alpha \in \mathbb{F}_{2^{n}}$ and $\beta \in \mathbb{F}_{2^{n}} \backslash\{0\}$ such that

$$
\left|\mathcal{W}_{F}(\alpha, \beta)\right| \geq 2^{(n+1) / 2}
$$

see [29], for instance. If $n$ is odd, there are functions $F$ with

$$
\left|\mathcal{W}_{F}(\alpha, \beta)\right| \leq 2^{(n+1) / 2}
$$

for all $\beta \neq 0$. Functions $F: \mathbb{F}_{2}^{n} \rightarrow \mathbb{F}_{2}^{n}$ with $\left|\mathcal{W}_{F}(\alpha, \beta)\right| \leq 2^{(n+1) / 2}$ for all $\beta \neq 0$ are called almost bent (AB). Note that AB functions can exist only if $n$ is odd. It is well known that any almost bent function is also APN (see [20]), but not vice versa, see the comments below. However, any quadratic APN function (see Definition 2) in $\mathbb{F}_{2}{ }^{n}$ must be AB , see [19]. If a function $F$ with $F(0)=0$ is AB , its Walsh spectrum is completely known:


\begin{equation*}
\left\{* \quad 2^{n}[1], 0\left[\left(2^{n-1}+1\right)\left(2^{n}-1\right)\right], \pm 2^{(n+1) / 2}\left[\left(2^{n}-1\right)\left(2^{n-2} \pm 2^{(n-3) / 2}\right)\right] *\right\} \tag{2}
\end{equation*}


(the values in brackets [ ] denote the multiplicities of the Walsh coefficients, and the notation $\{* \quad *\}$ indicates multisets). Similarly, the Walsh spectra of the Gold APN's (see Table 1) with $n$ even are completely known, too, see [21], for instance:


\begin{gather*}
\left\{* \quad 2^{n}[1], 0\left[\left(2^{n}-1\right)\left(2^{n-2}+1\right)\right], \pm 2^{(n+2) / 2}\left[\frac{1}{3}\left(2^{n}-1\right)\left(2^{n-3} \pm 2^{(n-4) / 2}\right)\right]\right.  \tag{3}\\
\left. \pm 2^{n / 2}\left[\frac{2}{3}\left(2^{n}-1\right)\left(2^{n-1} \pm 2^{(n-2) / 2}\right)\right] *\right\}
\end{gather*}


We say that an APN function with $F(0)=0$ and spectrum (2) (if $n$ is odd) or (3) (if $n$ is even) has the classical Walsh spectrum. We want to emphasize that the APN property alone does not determine the Walsh spectrum. APN functions may have quite different Walsh spectra. The reader can find the classical spectra in [21], for instance. If $F(0) \neq 0$, the distribution of the spectral values may be different, however the distribution of the absolute values does not change, see the comments following Proposition 1.

\begin{table}[h]
\begin{center}
\begin{tabular}{|l|l|l|l|}
\hline
 & Exponents d & Conditions & Proven in \\
\hline
Gold functions & $2^{i}+1$ & $\operatorname{gcd}(i, n)=1$ & [28, 36] \\
\hline
Kasami functions & $2^{2 i}-2^{i}+1$ & $\operatorname{gcd}(i, n)=1$ & [31, 32] \\
\hline
Welch function & $2^{t}+3$ & $n=2 t+1$ & [25] \\
\hline
Niho function & \( \begin{gathered} 2^{t}+2^{\frac{t}{2}}-1, t \text { even } \\ 2^{t}+2^{\frac{3 t+1}{2}}-1, t \text { odd } \end{gathered} \) & $n=2 t+1$ & [24] \\
\hline
Inverse function & $2^{2 t}-1$ & $n=2 t+1$ & [36] \\
\hline
Dobbertin function & $2^{4 t}+2^{3 t}+2^{2 t}+2^{t}-1$ & $n=5 t$ & [26] \\
\hline
\end{tabular}
\captionsetup{labelformat=empty}
\caption{Table 1: Known APN power functions $x^{d}$ on $\mathbb{F}_{2^{n}}$}
\end{center}
\end{table}

Throughout this paper, we always use the identification of the additive group of the vector space $\mathbb{F}_{2}{ }^{n}$ with the additive group of the finite field $\mathbb{F}_{2^{n}}$. Most APN and AB functions are defined as polynomial functions over $\mathbb{F}_{2^{n}}$, hence the constructions use the multiplicative structure of $\mathbb{F}_{2^{n}}$. However, the APN and AB property is just a property of the additive group of $\mathbb{F}_{2^{n}}$.

In Table 1, we list all known power APN mappings on $\mathbb{F}_{2^{n}}$ which are known so far: The Welch and Niho functions are also AB , the Gold and Kasami functions are AB if $n$ is odd. It is known that the inverse function and the Dobbertin function are not AB : The Walsh coefficients of the inverse function have been determined in [34], those of the Dobbertin function in [18], and these do not coincide with the values in (2) and (3). We note that the complete Walsh spectrum, i.e. the Walsh coefficients together with their multiplicities, are still not known for these two functions.

There are two questions which arise quite naturally:\\
Problem 1. (1.) Is the list of APN power functions in Table 1 complete?\\
(2.) Are all these examples "different"?

We will discuss these questions in the next section.

\section*{2 Equivalence of APN mappings}
Let us begin with the second part of the problem above. In order to describe whether two functions $F$ and $H$ are equivalent, we introduce group ring notation. This notion is also quite useful to describe the technique of "switching" an APN function. This is a very powerful tool to construct new APN functions, as we will show in this paper.

Let $\mathbb{F}$ be an arbitrary field, and let $(G,+)$ be an additively written abelian group (we are only interested in abelian groups, so we do not care about the general case). The group algebra $\mathbb{F}[G]$ consists of all "formal" sums

$$
\sum_{g \in G} a_{g} g, \quad a_{g} \in \mathbb{F} .
$$

We define componentwise addition

$$
\sum_{g \in G} a_{g} g+\sum_{g \in G} b_{g} g:=\sum_{g \in G}\left(a_{g}+b_{g}\right) g,
$$

and multiplication by

$$
\sum_{g \in G} a_{g} g \cdot \sum_{g \in G} b_{g} g:=\sum_{g \in G}\left(\sum_{h \in G} a_{h} \cdot b_{g-h}\right) g .
$$

Together with these two operations and the scalar multiplication $\lambda \sum_{g \in G} a_{g} g:=\sum_{g \in G}\left(\lambda a_{g}\right) g$, the set $\mathbb{F}[G]$ becomes an algebra, the so called group algebra. The dimension of this algebra as an\\
$\mathbb{F}$-vectorspace is $|G|$. Given a function $F: \mathbb{F}_{2}^{n} \rightarrow \mathbb{F}_{2}^{n}$, we associate a group algebra element $G_{F}$ in $\mathbb{F}\left[\mathbb{F}_{2}{ }^{n} \times \mathbb{F}_{2}{ }^{n}\right]$ with it:

$$
G_{F}:=\sum_{v \in \mathbb{F}_{2}^{n}}(v, F(v)) \quad \text { in } \mathbb{F}\left[\mathbb{F}_{2}^{n} \times \mathbb{F}_{2}^{n}\right]
$$

The coefficients of the group elements in $G_{F}$ are just 0 or 1 . More generally, any subset $T$ of a group $G$ can be identified with the element $\sum_{g \in T} g$, where the coefficients of all elements in $T$ are 1 , and the coefficients of elements not in $T$ are 0 . Hence, $G_{F}$ is the group algebra element corresponding to the "graph" of the function $F$, which consists of all pairs $(v, F(v)), v \in \mathbb{F}_{2}{ }^{n}$.

We have the following Lemma:\\
Lemma 1. A function $F: \mathbb{F}_{2}^{n} \rightarrow \mathbb{F}_{2}^{n}$ is $A P N$ if and only if


\begin{equation*}
G_{F} \cdot G_{F}=2^{n} \cdot(0,0)+2 \cdot D_{F} \text { in } \mathbb{C}\left[\mathbb{F}_{2}^{n} \times \mathbb{F}_{2}^{n}\right] \tag{4}
\end{equation*}


for some subset $D_{F} \in\left(\mathbb{F}_{2}{ }^{n} \times \mathbb{F}_{2}{ }^{n}\right) \backslash\{(0,0)\}$.\\
Proof. The coefficient of ( $u, v$ ) in $G_{F} \cdot G_{F}$ is the number of pairs ( $x, F(x)$ ) and ( $y, F(y)$ ) such that $(x+y, F(x)+F(y))=(u, v)$. This is the number of solutions $F(x)+F(y)=v$ where $x+y=u$, hence it is the number of solutions of $F(x)+F(x+u)=v$. This integer is $2^{n}$ for $(u, v)=(0,0)$, otherwise it is 0 or 2 (since $F$ is APN).

In (4), we may replace $\mathbb{C}$ by any field of characteristic $\neq 2$ if we add for some subset $D_{F} \in \left(\mathbb{F}_{2}^{n} \times \mathbb{F}_{2}^{n}\right) \backslash\{(0,0)\}$ of size $2^{n-1} \cdot\left(2^{n}-1\right)$.

We emphasize that $G$ is additively written, but this addition is quite different from the addition in the group algebra $\mathbb{F}[G]$. If, for instance, $A, B \subset G$ and $A \cap B=\emptyset$, then $A \cup B$ is the subset of $G$ corresponding to $A+B$ in $\mathbb{F}[G]$. If $g \in G$, then $A \cdot g$ in $\mathbb{F}[G]$ corresponds to the subset $\{a+g: a \in A\}$. We call $A \cdot g$ a translate of $A$. It looks a bit awkward that the product $A \cdot g$ is the set of sums $a+g$ with $a \in A$.

The ideal generated by $G_{F}$ in $\mathbb{F}_{2}\left[\mathbb{F}_{2}^{n} \times \mathbb{F}_{2}^{n}\right]$ is a subspace of the $2^{2 n}$-dimensional vector space of the group algebra $\mathbb{F}_{2}\left[\mathbb{F}_{2}{ }^{n} \times \mathbb{F}_{2}{ }^{n}\right]$. The dimension is called the $\Gamma$-rank of the function $F$. Similarly, the dimension of the ideal generated by $D_{F}$ in $\mathbb{F}_{2}\left[\mathbb{F}_{2}{ }^{n} \times \mathbb{F}_{2}{ }^{n}\right]$ is called the $\Delta$ - $\mathbf{r a n k}$ of $F$.

The Walsh transform of a function $F$ is nothing else than the Discrete Fourier transform of $G_{F}$, which we describe briefly: If $G$ is a finite abelian group, then there are $|G|$ different homomorphisms $\chi: G \rightarrow \mathbb{C}$, and the set of these homomorphisms (called characters) form a group under multiplication $\chi_{1} \chi_{2}(g):=\chi_{1}(g) \cdot \chi_{2}(g)$. This group is isomorphic to $G$. Characters $\chi$ may be extended to homomorphisms $\chi: \mathbb{C}[G] \rightarrow \mathbb{C}[G]$ by linearity:

$$
\chi\left(\sum_{g \in G} a_{g} g\right):=\sum_{g \in G} a_{g} \chi(g)
$$

Let $\chi$ be a character of $G$, and $\Psi$ an automorphism of $G$. Then the mapping $\chi^{\Psi}: G \rightarrow \mathbb{C}$ with $\chi^{\Psi}(g):=\chi(\Psi(g))$ is again a character. Moreover, $\Psi$ may be extended to a group algebra automorphism. This shows that the Walsh spectrum of an element $D \in \mathbb{C}[G]$ is invariant under the application of group automorphisms.

If $G=\mathbb{F}_{2}^{n} \times \mathbb{F}_{2}^{n}$, the characters are the mappings $\chi_{\alpha, \beta}$ defined by $\chi_{\alpha, \beta}(u, v):=(-1)^{\operatorname{tr}(\alpha u+\beta v)}$, where we identify (as usual) the vector space $\mathbb{F}_{2}^{n}$ with the additive group of the finite field $\mathbb{F}_{2^{n}}$. Therefore, the Walsh spectrum is just the multiset of character values of $G_{F}$.

Definition 1 (CCZ and EA equivalence, [14]). Two functions $F, H: \mathbb{F}_{2}{ }^{n} \rightarrow \mathbb{F}_{2}{ }^{n}$ are called CCZ equivalent if there is an automorphism $\Psi$ of the group $\mathbb{F}_{2}{ }^{n} \times \mathbb{F}_{2}{ }^{n}$ (which is the additive group of the vector space $\mathbb{F}_{2}{ }^{2 n}$, hence $\Psi$ is simply a bijective linear mapping on this vector space) and an element $(u, v) \in \mathbb{F}_{2}{ }^{n} \times \mathbb{F}_{2}{ }^{n}$ such that

$$
\Psi\left(G_{F}\right)=G_{H} \cdot(u, v)
$$

hence $\Psi\left(G_{F}\right)$ is a translate of $G_{H}$. If $\Psi$ fixes the subgroup $\left\{(0, y): y \in \mathbb{F}_{2}^{n}\right\}$ setwise, we say that the functions are $\mathbf{E A}$ equivalent $(E A=$ extended affine $)$. If, additionally, $\Psi$ fixes the set $\left\{(x, 0): x \in \mathbb{F}_{2}{ }^{n}\right\}$, then $F$ and $G$ are called affine equivalent.

This relation is called $C C Z$ equivalent because it has been introduced (using different notation) by Carlet, Charpin and Zinoviev [19].

Proposition 1. If $F$ is an $A P N$ (resp. $A B$ ) function, and if $H$ is $C C Z$ equivalent to $F$, then $H$ is also an $A P N$ (resp. $A B$ ) function.

Proof. Let $G_{H}=\Psi\left(G_{F}\right) \cdot(u, v)$. If $\chi$ is a character of $\mathbb{F}_{2}{ }^{n} \times \mathbb{F}_{2}{ }^{n}$, then $\chi\left(G_{H}\right)=\chi^{\Psi}\left(G_{F}\right) \cdot \chi(u, v)$, hence the maximum absolute character value is invariant under CCZ equivalence, hence $H$ is AB if $F$ is $A B$. If $F$ is APN, then

$$
G_{F} \cdot G_{F}=n \cdot(0,0)+2 \cdot D_{F}
$$

and therefore

$$
G_{H} \cdot G_{H}=n \cdot(0,0)+2 \cdot \Psi\left(D_{F}\right) \quad \text { in } \mathbb{C}\left[\mathbb{F}_{2}^{n} \times \mathbb{F}_{2}^{n}\right]
$$

This proposition and its proof have some consequences: The Walsh spectrum is not invariant under CCZ equivalence: The Walsh coefficients $\chi\left(G_{F}\right)$ and $\chi\left(G_{F}\right) \cdot \chi(u, v)$ differ by the factor $\chi(u, v)$, hence by $\pm 1$. The problem arises from the addition of the element $(u, v)$ : The Walsh spectrum is invariant under affine equivalence, but not under EA or CCZ equivalence. The set containing the Walsh spectrum and its negative is called the extended Walsh spectrum, and this is invariant under CCZ equivalence. Sometimes we just speak about the Walsh spectrum, for instance when discussing the tables at the end of this paper. In these cases, we list just one of the two possible Walsh spectra.

There is one drawback in the concept of CCZ equivalence: If $F$ is APN (or AB ), the group algebra element $\Psi\left(G_{F}\right)$ does not necessarily correspond to a function $H$, see [14], for instance.

The question about the inequivalence of the power APN functions in Table 1 has been discussed in [12].

It is obvious that the $\Delta$ - and $\Gamma$-ranks are invariant under CCZ equivalence.\\[0pt]
Now we discuss the first part of the problem at the end of Section 1: Is the list in Table 1 complete? This has been answered negatively in [27]. One of the examples in [27] has been generalized to an infinite family, and a lot more constructions have been found since. In particular, Dillon [22] presented a list of 12 examples in $\mathbb{F}_{2}^{6}$. This list appears in [9], together with many new examples in the cases $\mathbb{F}_{2}^{7}$ and $\mathbb{F}_{2}^{8}$.

However, all the new examples that have been constructed so far are "quadratic" in the sense that the derivatives $F(x+a)-F(x)$ are affine mappings. Since the property of being "quadratic" is not invariant under CCZ equivalence (see [14]), we modify the definition as follows:

Definition 2. A function $F: \mathbb{F}_{2}{ }^{n} \rightarrow \mathbb{F}_{2}{ }^{n}$ is $\mathbf{C C Z}$ quadratic if $F$ is $C C Z$ equivalent to a function $H$ with the property that $H(x+a)-H(x)$ is affine for all $a \in \mathbb{F}_{2}{ }^{n}$.

How can we prove that a function is not CCZ quadratic? For this purpose, we look at the following design or incidence structure associated with an APN function. We refer the reader to the comprehensive book [2] for background from design theory and difference sets: The designs we are going to define here may be viewed as the designs developed from a certain type of difference set.

Definition 3 ([30]). Let $F: \mathbb{F}_{2}{ }^{n} \rightarrow \mathbb{F}_{2}{ }^{n}$ be an APN function. Then we define two incidence structures (designs) on the point set $\mathbb{F}_{2}{ }^{n} \times \mathbb{F}_{2}{ }^{n}$ : The first design has blocks

$$
G_{F} \cdot(a, b):=\left\{(x+a, F(x)+b): x \in \mathbb{F}_{2}^{n}\right\}
$$

for $a, b \in \mathbb{F}_{2}^{n}$, i.e. the translates of $G_{F}$. We call this design the development of $G_{F}$, denoted by $\operatorname{dev}\left(G_{F}\right)$. Similarly, the design whose blocks are the translates

$$
D_{F} \cdot(a, b)
$$

of $D_{F}$ (see Lemma 1) is the development of $D_{F}$, denoted by $\operatorname{dev}\left(D_{F}\right)$. We call two designs isomorphic if there is a permutation $\pi$ on the set of points such that blocks (which in this case are subsets $B=\left\{g_{1}, \ldots\right\}$ of the point set) are mapped to blocks (i.e. $\pi(B)=\left\{\pi\left(g_{1}\right), \ldots\right\}$ is a block).

Any incidence structure gives rise to an incidence matrix: Rows and columns are indexed by the points and blocks, resp., and the ( $p, B$ )-entry is 1 if point $p$ is incident with block $B$; all other entries are 0 . The $\Gamma$-rank defined earlier is nothing else than the rank of the incidence matrix of $\operatorname{dev}\left(G_{F}\right)$, considered as a matrix with entries in $\mathbb{F}_{2}$; similarly, the $\Delta$-rank is the $\mathbb{F}_{2}$-rank of an incidence matrix of $\operatorname{dev}\left(D_{F}\right)$.

Lemma 2. If $F$ and $H$ are $C C Z$ equivalent $A P N$ functions, the $\operatorname{designs} \operatorname{dev}\left(G_{F}\right)$ and $\operatorname{dev}\left(G_{H}\right)$ are isomorphic. Moreover, the $\operatorname{designs} \operatorname{dev}\left(D_{F}\right)$ and $\operatorname{dev}\left(D_{H}\right)$ are isomorphic.

Proof. Straightforward, see also [30]: The group automorphism $\Psi$ with $\Psi\left(G_{F}\right)=G_{H} \cdot(u, v)$ is the permutation on the point set which maps blocks to blocks.

Using MAGMA [4] it is quite easy to determine the automorphism groups of these designs for small values of $n$. There is another group associated with the $\operatorname{designs} \operatorname{dev}\left(G_{F}\right)$ (resp. $\operatorname{dev}\left(D_{F}\right)$ ): The sets $G_{F}$ (resp. $D_{F}$ ) are subsets of $\mathbb{F}_{2}{ }^{2 n}$. Then there may exist automorphisms $\varphi$ of $\mathbb{F}_{2}{ }^{2 n}$ such that $\varphi\left(G_{F}\right)=G_{F} \cdot(u, v)$ (resp. $\varphi\left(D_{F}\right)=D_{F} \cdot(u, v)$ ) for some $u, v \in \mathbb{F}_{2}{ }^{n}$. These automorphisms form a group contained in the automorphism group of the designs $\operatorname{dev}\left(G_{F}\right)$ (resp. $\operatorname{dev}\left(D_{F}\right)$ ). Using notation adopted from the theory of difference sets, we call the group of these automorphisms the multiplier group $\mathcal{M}\left(G_{F}\right)$ (resp. $\mathcal{M}\left(D_{F}\right)$ ) of $\operatorname{dev}\left(G_{F}\right)$ (resp. $\operatorname{dev}\left(D_{F}\right)$ ). It turns out that $\mathcal{M}\left(G_{F}\right)$ is much easier to compute with MAGMA than the full automorphism group of the design $\operatorname{dev}\left(G_{F}\right)$.

We denote the group of translations $\tau_{a, b}: \mathbb{F}_{2}^{n} \times \mathbb{F}_{2}^{n} \rightarrow \mathbb{F}_{2}^{n} \times \mathbb{F}_{2}^{n}$ with $\tau_{a, b}(x, y):=(x+a, y+b)$ by $\mathcal{T}$. Obviously, we have $|\mathcal{T}|=2^{2 n}$, and $\left|\mathcal{M}\left(G_{F}\right) \cap \mathcal{T}\right|=1$ as well as $\left|\mathcal{M}\left(D_{F}\right) \cap \mathcal{T}\right|=1$.

Since the multiplier group normalizes $\mathcal{T}$, we have the following lemma:\\
Lemma 3. Let $F: \mathbb{F}_{2}{ }^{n} \rightarrow \mathbb{F}_{2}{ }^{n}$ be an $A P N$ function. Then

\begin{enumerate}
  \item $\left\langle\mathcal{M}\left(G_{F}\right), \mathcal{T}\right\rangle \subseteq A u t\left(\operatorname{dev}\left(G_{F}\right)\right)$.
  \item $\left\langle\mathcal{M}\left(D_{F}\right), \mathcal{T}\right\rangle \subseteq \operatorname{Aut}\left(\operatorname{dev}\left(D_{F}\right)\right)$.
  \item $\left|\mathcal{M}\left(G_{F}\right)\right| \cdot 2^{2 n}=\left|\left\langle\mathcal{M}\left(G_{F}\right), \mathcal{T}\right\rangle\right|$.
  \item $\left|\mathcal{M}\left(D_{F}\right)\right| \cdot 2^{2 n}=\left|\left\langle\mathcal{M}\left(D_{F}\right), \mathcal{T}\right\rangle\right|$.
  \item $\mathcal{M}\left(G_{F}\right) \subseteq \mathcal{M}\left(D_{F}\right)$.
\end{enumerate}

It is possible to show that $\mathcal{M}\left(G_{F}\right)$ is just the automorphism group of the extended code $\mathcal{C}_{F}^{\text {ext }}$ defined in the introduction, see [9]. In all cases known to us, the "full" automorphism group of the $\operatorname{design} \operatorname{dev}\left(G_{F}\right)$ is just the multiplier group "times" the translations $\tau_{a, b}$.

We do not know whether this observation that holds for small values is true in general:\\
Problem 2. Is it possible that the full automorphism group of $\operatorname{dev}\left(G_{F}\right)$ (resp. $\operatorname{dev}\left(D_{F}\right)$ ) is larger than $\left|\mathcal{M}\left(G_{F}\right)\right| \cdot 2^{2 n}\left(\right.$ resp.$\left.\left|\mathcal{M}\left(D_{F}\right)\right| \cdot 2^{2 n}\right)$ ?

It seems that the automorphism group is a good invariant for CCZ equivalence, in particular to distinguish the quadratic from the non-quadratic case:

Theorem 4. If $F: \mathbb{F}_{2}^{n} \rightarrow \mathbb{F}_{2}^{n}$ is an APN mapping such that $F(x+a)-F(x)$ is affine for all $a \in \mathbb{F}_{2}{ }^{n}$, then the multiplier group $\mathcal{M}\left(G_{F}\right)$ contains an elementary abelian group of order $2^{n}$.

Proof. If $F$ is quadratic, then we may assume (after replacing $F$ by a CCZ equivalent function, if necessary) that the mappings $L_{a}(x):=F(x+a)+F(x)+F(a)+F(0)$ are linear. We compute

$$
\begin{aligned}
\left(L_{a}+L_{b}\right)(x) & =F(x)+F(x+a)+F(a)+F(0)+F(x)+F(x+b)+F(b)+F(0) \\
& =F(x+a)+F(x+b)+F(a)+F(b) .
\end{aligned}
$$

Now we use

$$
L_{a}(b+x)=F(a+b+x)+F(b+x)+F(a)+F(0)
$$

and

$$
L_{a}(b+x)=L_{a}(b)+L_{a}(x)=F(b)+F(a+b)+F(x)+F(a+x)
$$

to obtain


\begin{equation*}
F(a+b+x)=F(b+x)+F(a)+F(0)+F(b)+F(a+b)+F(x)+F(a+x) . \tag{5}
\end{equation*}


We get

$$
\begin{aligned}
L_{a+b}(x) & =F(a+b+x)+F(a+b)+F(x)+F(0) \\
& =F(b+x)+F(a)+F(b)+F(a+x) \text { using (5) } \\
& =L_{a}(x)+L_{b}(x)
\end{aligned}
$$

This shows that the mappings $\psi_{a}$ defined by $\psi_{a}(x, y)=\left(x, y+L_{a}(x)\right)$ are linear, and

$$
\begin{aligned}
\psi_{a}\left(G_{F}\right) & =\left\{\left(x, F(x)+L_{a}(x)\right): x \in \mathbb{F}_{2}^{n}\right\} \\
& =\left\{(x, F(x+a)+F(a)+F(0)): x \in \mathbb{F}_{2}^{n}\right\} \\
& =\left\{(x-a, F(x)+F(a)+F(0)): x \in \mathbb{F}_{2}^{n}\right\} \\
& =G_{F} \cdot(-a, F(a)+F(0))
\end{aligned}
$$

is a translate of $G_{F}$, hence the mappings $\psi_{a}$ are automorphisms of $\operatorname{dev}\left(G_{F}\right)$. Moreover, $\psi_{a+b}= \psi_{b} \circ \psi_{a}$ so that the $\psi_{a}$ 's form a group of order $2^{n}$.

With Lemma 3, we get the following\\
Corollary 5. Under the assumptions of Theorem 4, $\operatorname{Aut}\left(\operatorname{dev}\left(G_{F}\right)\right)$ and $\operatorname{Aut}\left(\operatorname{dev}\left(D_{F}\right)\right)$ both contain an elementary abelian group of order $2^{3 n}$.

Corollary 6 (Göloğlu, Pott [30]). The Kasami power functions $x^{13}$ and $x^{57}$ on $\mathbb{F}_{2^{7}}$ are not $C C Z$ quadratic, hence they are not CCZ equivalent to quadratic functions.

Proof. Using MAGMA, it is easy to compute $\left|\operatorname{Aut}\left(\operatorname{dev}\left(G_{F}\right)\right)\right|$ for $F(x)=x^{13}$ and $F(x)=x^{57}$ : The order of the groups is, in both cases, $2^{14} \cdot 7 \cdot\left(2^{7}-1\right)$ (Table 6) which is not divisible by $2^{21}$.

Most people conjecture that the examples in Table 1 are all CCZ inequivalent, except for small $n$ where some of the cases coincide, but as far as we know there is no proof yet. It is known that the Gold power mappings are CCZ inequivalent to the Kasami power mappings, and different Gold exponents are CCZ inequivalent, see [12].

There is another concept related to quadratic APN functions: If $F$ is quadratic, then $F(x+ a)-F(x)$ is affine and


\begin{equation*}
H_{a}:=\left\{b: F(x+a)-F(x)=b \quad \text { for some } x \in \mathbb{F}_{2}^{n}\right\} \tag{6}
\end{equation*}


is an affine subspace. If $a \neq 0$, this subspace has $2^{n-1}$ elements (since $F$ is APN), hence its is an (affine) hyperplane. We say that a function is crooked if the sets in (6) are (affine) hyperplanes for all $a \neq 0$. This concept is due to Bending and Fon-der-Flaas [1]. "Crooked" is not invariant under CCZ equivalence, hence it would be better to say that a function $F$ is crooked if it is CCZ equivalent to a function for which all the sets $H_{a}$ are hyperplanes:

Definition 4. An APN function $F: \mathbb{F}_{2}{ }^{n} \rightarrow \mathbb{F}_{2}{ }^{n}$ is called $\mathbf{C C Z}$ crooked if $F$ is $C C Z$ equivalent to a function $G$ such that the sets

$$
\left\{b: G(x+a)-G(x)=b \text { for some } x \in \mathbb{F}_{2}^{n}\right\}
$$

are affine hyperplanes in $\mathbb{F}_{2}{ }^{n}$.\\[0pt]
It is obvious that any (CCZ) quadratic function is (CCZ) crooked, and it is conjectured that the converse is also true, see [3,33] for partial results in this direction. However, as long as we do not know whether non-quadratic crooked functions exist, we need to find arguments that a function is not CCZ crooked. The following argument gives an interesting necessary condition for a function to be crooked:

Theorem 7. Let $F: \mathbb{F}_{2}^{n} \rightarrow \mathbb{F}_{2}^{n}$ be an APN mapping. If $F$ is $C C Z$ crooked, then the dimension of the ideal generated by $D_{F} \in \mathbb{F}_{2}\left[\mathbb{F}_{2}^{n} \times \mathbb{F}_{2}^{n}\right]$ is at most $2^{n+1}$, hence the $\Delta$-rank is at most $2^{n+1}$ (see Lemma 1 for the definition of $D_{F}$ ).

Proof. If $F$ is crooked, there are $2^{n}-1$ (affine) hyperplanes $H_{a}$ such that

$$
D_{F}=\left\{(a, x): a \in \mathbb{F}_{2}^{n} \backslash\{0\}, x \in H_{a}\right\}
$$

(replace $F$ by a CCZ equivalent function if necessary). We define

$$
J_{a}:=\left\{(a, x): x \in \mathbb{F}_{2}^{n}\right\}
$$

As explained above, these subsets may be also interpreted as elements in $\mathbb{F}_{2}\left[\mathbb{F}_{2}^{n} \times \mathbb{F}_{2}^{n}\right]$. We will show that the ideal generated by $D_{F}$ is contained in the subspace $\mathcal{I}$ generated (as a vector space) by the $2^{n+1}$ elements

$$
\left\{D_{F} \cdot(u, 0): u \in \mathbb{F}_{2}^{n}\right\} \cup\left\{J_{a}: a \in \mathbb{F}_{2}^{n}\right\}
$$

in the group algebra $\mathbb{F}_{2}\left[\mathbb{F}_{2}^{n} \times \mathbb{F}_{2}^{n}\right]$. It is sufficient to show that $D_{F} \cdot(u, v) \in \mathcal{I}$ for all $(u, v) \in \mathbb{F}_{2}{ }^{n} \times \mathbb{F}_{2}{ }^{n}$. The set corresponding to $D_{F} \cdot(u, v)$ is $\left\{\left(a+u, H_{a}+v\right): a \in \mathbb{F}_{2}{ }^{n} \backslash\{0\}\right\}$, where $H_{a}+v=\left\{h+v: h \in H_{a}\right\}$. Here we used the notation ( $x, T$ ) to denote the set of elements $\{(x, t): t \in T\}$. Since $H_{a}$ is a hyperplane, we have $H_{a}+y=H_{a}$ or $H_{a}+y$ is the complement of $H_{a}$ : In group algebra notation, this means for fixed $a \in \mathbb{F}_{2}^{n} \backslash\{0\}$

$$
\left(a+u, H_{a}+v\right)=\left(a+u, H_{a}\right) \quad \text { or } \quad\left(a+u, H_{a}+v\right)=\left(a+u, H_{a}\right)+J_{a+u}
$$

in $\mathbb{F}_{2}\left[\mathbb{F}_{2}^{n} \times \mathbb{F}_{2}^{n}\right]$. In the equation above, we again identify subsets with the corresponding group algebra elements. Adding the element $J_{a+u}$ has the effect of complementing $H_{a}$ in ( $a+u, H_{a}$ ).

Corollary 8. The Kasami power mappings $x^{13}$ and $x^{57}$ on $\mathbb{F}_{2^{7}}$ are not $C C Z$ crooked.\\
Proof. It is easy to compute the $\Delta$-ranks of $x^{13}$ (resp. $x^{57}$ ) using MAGMA: The ranks are 338 (resp. 436), see Table 6.

It would be very interesting to determine the $\Delta$ - and $\Gamma$-ranks of APN functions theoretically.\\
In the next section (Theorem 11), we will construct a new APN function which, at first glance, seems to be non-quadratic. In order to prove that the function is indeed non-quadratic, we use Theorem 7 to show that the function cannot be CCZ equivalent to a crooked function, hence it cannot be quadratic. We could also use Theorem 4 to show that the function is non-quadratic, since the automorphism group of $\operatorname{dev}\left(G_{F}\right)$ is too small for the new function $F$. We have checked that our function is equivalent to the new example given in [7]. However, in that paper the authors erroneously claimed that their new function is CCZ equivalent to a quadratic one. Moreover, our function was found independently by the search in [7].

\section*{3 Dillon's switching construction}
The following interesting construction of an APN function is contained in [11]:\\
Proposition 2. The function $x^{3}+\operatorname{tr}\left(x^{9}\right)$ is APN in $\mathbb{F}_{2^{n}}$ for all $n$.\\
This is a special case of what we call "switching". For this purpose, we consider certain projection homomorphisms on the group algebra $\mathbb{F}[G]$. Let $U$ be a subgroup of $G$. Then the canonical homomorphism $\varphi_{U}: G \rightarrow G / U$ defined by $\varphi_{U}(g):=g+U$ can be extended by linearity to a homomorphism $\varphi_{U}: \mathbb{F}[G] \rightarrow \mathbb{F}[G / U]$. Let $D=\sum a_{g} g$ be an element in $\mathbb{F}[G]$. The coefficient of $g+U$ in $\varphi(D)$ is $\sum_{h \in g+U} a_{h}$ (this summation is in $\mathbb{F}$ ). If $D$ has only coefficients 0 and 1 (so that $D$ corresponds to a set $D \subseteq G$ ) then the coefficient of $g+U$ is $|D \cap(g+U)|$. In particular, if each coset of $U$ meets $D$ in at most one element, then $\varphi_{U}(D)$ has also only coefficients 0 and 1 . This is the case if $U \leq\{0\} \times \mathbb{F}_{2}{ }^{n}$ (recall that " $\leq$ " denotes "being subgroup").

Definition 5 (switching neighbours). Let $F, H: \mathbb{F}_{2}{ }^{n} \rightarrow \mathbb{F}_{2}{ }^{n}$ be two functions, and let $U \leq \mathbb{F}_{2}{ }^{n} \times \mathbb{F}_{2}{ }^{n}$ be a subgroup of $\mathbb{F}_{2}^{n} \times \mathbb{F}_{2}^{n}$. We say that $F$ and $H$ switching neighbours with respect to $U$ if $\varphi_{U}\left(G_{F}\right)=\varphi_{U}\left(G_{H}\right)$. We call $F$ and $H$ switching neighbours in the narrow sense if $U \leq\{0\} \times \mathbb{F}_{2}{ }^{n}$ and $\operatorname{dim}(U)=1$.

If $F$ and $H$ are switching neighbours with respect to $U$, we may obtain $H$ from $F$ by first projecting $G_{F}$ onto $\varphi_{U}\left(G_{F}\right)$, and then lifting this element to $G_{H}$. We may also try to construct new switching neighbours $H$ of $F$ via such a project and lift procedure such that (hopefully) $F$ and $H$ are CCZ inequivalent. This is promising in particular if the dimension of $U$ is small. The intuitive idea behind this approach is that $\varphi_{U}\left(G_{F}\right)$ is almost an APN function, and so it may be easy to turn this "almost" APN into an APN function.

We describe this approach (and applications) for the case $F: \mathbb{F}_{2}{ }^{n} \rightarrow \mathbb{F}_{2}{ }^{n}$ and $U \leq\{0\} \times \mathbb{F}_{2}{ }^{n}$. This has the advantage that the coefficients of $\varphi_{U}\left(G_{F}\right)$ are 0 and 1 only, since the cosets of $\{0\} \times \mathbb{F}_{2}^{n}$ (and therefore also the cosets of $U$ ) meet $G_{F}$ no more than once. In this case, $\varphi_{U}\left(G_{F}\right)$ corresponds to a mapping $F_{U}: \mathbb{F}_{2}{ }^{n} \rightarrow \mathbb{F}_{2}{ }^{n} / U^{\prime}$ with $F_{U}(v):=v+U^{\prime}$ and


\begin{equation*}
U^{\prime}=\{u:(0, u) \in U\} \tag{7}
\end{equation*}


(hence $U^{\prime}$ is basically the same as $U$ ).\\
Proposition 3. Let $F, H: \mathbb{F}_{2}^{n} \rightarrow \mathbb{F}_{2}^{n}$, and let $U \leq\{0\} \times \mathbb{F}_{2}^{n}$. Then

$$
F_{U}=H_{U} \quad \text { if and only if } \quad(0, F(v)-H(v)) \in U \quad \text { for all } v \in \mathbb{F}_{2}^{n} .
$$

If $U=\{(0,0),(0, u)\}$, then $F_{U}=H_{U}$ if and only if there is a Boolean function $f: \mathbb{F}_{2}^{n} \rightarrow \mathbb{F}_{2}$ such that $H(v)=F(v)+f(v) \cdot u$.

Proof. We define $U^{\prime}$ as in (7). Then $F_{U}(v)=H_{U}(v)$ if and only if $F(v)+U^{\prime}=H(v)+U^{\prime}$, hence $F_{U}(v)-H_{U}(v) \in U^{\prime}$ for all $v$. This shows the first part of the proposition.

The function $f$ is defined via

$$
f(v):= \begin{cases}0 & \text { if } F(v)=H(v) \\ 1 & \text { if } F(v) \neq H(v)\end{cases}
$$

which concludes the proof.\\
The two functions $F(x)=x^{3}$ and $H(x)=x^{3}+\operatorname{tr}\left(x^{9}\right)$ are switching neighbours in the narrow sense: Take the 1 -dimensional subspace $U$ generated by $(0,1) \in \mathbb{F}_{2}^{n} \times \mathbb{F}_{2}^{n}$. Then $\varphi_{U}\left(G_{F}\right)= \varphi_{U}\left(G_{H}\right)$.

Proposition 3 shows that we may obtain all switching neighbours of $F$ in the narrow sense (with respect to a one-dimensional subspace) by adding a Boolean function $f$ times a vector $u \neq 0$. Let $F$ be an APN function. The following Theorem gives a necessary and sufficient condition for $f$ to produce another (not necessarily equivalent) APN function:

Theorem 9. Assume that $F: \mathbb{F}_{2}{ }^{n} \rightarrow \mathbb{F}_{2}{ }^{n}$ is an APN function. Let $u \in \mathbb{F}_{2}{ }^{n}, u \neq 0$, and let $f: \mathbb{F}_{2}{ }^{n} \rightarrow \mathbb{F}_{2}$ be a Boolean function. Then $F(v)+f(v) \cdot u$ is an APN function if and only if

$$
f(x)+f(x+a)+f(y)+f(y+a)=0
$$

for all $x, y, a \in \mathbb{F}_{2}{ }^{n}$ with


\begin{equation*}
F(x)+F(x+a)+F(y)+F(y+a)=u \tag{8}
\end{equation*}


Proof. Since $F$ is APN, the equation

$$
F(x+a)+F(x)+(f(x+a)+f(x)) u=b
$$

hat at most 4 solutions for $x$, namely those $x$ for which $F(x+a)+F(x) \in\{b, b+u\}$. If there are 4 different solutions $x, y, x+a, y+a$ (hence if $F$ is not APN) then

$$
\begin{aligned}
F(x+a)+F(x)+(f(x+a)+f(x)) u & =b \\
F(y+a)+F(y)+(f(y+a)+f(y)) u & =b .
\end{aligned}
$$

But this is possible if and only if (8) and

$$
f(x)+f(x+a)+f(y)+f(y+a)=1
$$

holds. Therefore, $F$ is APN if $f(x)+f(x+a)+f(y)+f(y+a)=0$.\\
Remark 1. 1. The Boolean function $f$ depends on $u$, i.e. for different choices of $u$ we may get different $f$ 's.\\
2. It seems to be difficult to find a theoretical criterion for the function $F(v)+f(v) u$ to be $C C Z$ equivalent to $F(v)$.\\
3. The functions $F(v)$ and $F(v)+f(v) u$ are switching neighbours in the narrow sense with respect to $\{(0,0),(0, u)\}$.

Theorem 9 immediately suggests a strategy to find Boolean functions $f$ such that $F(v)+f(v) u$ is APN: Determine all 4 -tuples $x, y, x+y, y+a$ such that ( 8 ) holds. These 4 -tuples give rise to constraints

$$
f(x)+f(x+a)+f(y)+f(y+a)=0
$$

We may view $f$ as a vector of length $2^{n}$ (coordinates are indexed by elements $v$ in $\mathbb{F}_{2}{ }^{n}$, and the entries of the vector are $f(v)$ ). Thus the constraints are linear conditions, and we may find $f$ 's by solving a system of linear equations.

Here is another interpretation:\\
The function $F$ is uniquely determined by its function values. Consider the $n$-dimensional subspace $W$ of $\mathbb{F}_{2}^{2^{n}}$ spanned by the rows of $\left(F(x): x \in \mathbb{F}_{2}{ }^{n}\right)$. Write $\mathbb{F}_{2}{ }^{n}$ as a direct sum $u \mathbb{F}_{2} \oplus \bar{U}$. This "lifts" to a decomposition of $W=V \oplus \bar{V}$ : We simply decompose every $F(x) \in \mathbb{F}_{2}^{n}$ according to a decomposition $\mathbb{F}_{2}{ }^{n}=\langle u\rangle \oplus \bar{U}$, where $\bar{U}$ is any complement of $\langle u\rangle$.

Let $x, y, x+a, y+a$ be a 4-tuple which satisfies the condition of Theorem 9. The corresponding indicator function (which is 1 for $x, y, x+a, y+a$ and 0 otherwise) is a vector of weight 4 in $\bar{V}^{\perp}$. This function is not in $W^{\perp}$ since $F$ is APN . Let $\tilde{V}$ be the vector space generated by these vectors of weight 4 . Note that $F(x)+f(x) \cdot u$ is APN if and only if $f \in \tilde{V}^{\perp} 1$ (Theorem 9). Let $R$ be the vector space of dimension $n+1$ generated by the all-one-vector $(1, \ldots, 1)$ and the rows of ( $x: x \in \mathbb{F}_{2}^{n}$ ). Then we have $\bar{V} \oplus R \subseteq \tilde{V}^{\perp}$. If $f, g \in \tilde{V}^{\perp}$ are in the same coset of $\bar{V} \oplus R$, then $F(x)+f(x) u$ and $F(x)+g(x) u$ are EA equivalent: The difference $\left(f(x)-g(x): x \in \mathbb{F}_{2}^{n}\right)$ is an element of $\bar{V} \oplus R \subseteq W \oplus R$, hence $f(x)-g(x)=A(F(x))+B(x)$ or $A(F(x))+B(x)+1$ for some linear mappings $A, B_{\tilde{\nu}}: \mathbb{F}_{2}{ }^{n} \mapsto \mathbb{F}_{2}$.

In particular: If $\tilde{V}^{\perp}$ has dimension $2 n$, then there is no candidate for a "switching function" $f(x) u$.

\footnotetext{${ }^{1}$ the ${ }^{\perp}$ is missing in the printed article
}Definition 6. The finest equivalence relation on the set of APN functions such that all switching neighbours in the narrow sense and all $E A$-equivalent functions are equivalent, is called the $\mathbf{E} \mathbf{A}$ switching equivalence relation. In the same way, we define CCZ switching equivalence.

This "switching idea" is closely related to a comment of John F. Dillon about the construction in [11]. Actually our proof of Theorem 9 is just a slight variation of Dillon's observation: Using the notation of Proposition 3, he considered the case $u=1$. Therefore we decided to call this method Dillon's switching construction since he was the first who discussed it.

In the next section, we determine the CCZ switching equivalence classes of all known CCZ inequivalent APN functions on $\mathbb{F}_{2}{ }^{n}, n \leq 7$, and all EA switching equivalence classes of all known CCZ inequivalent APN functions if $n=8$ and $n=9$. Several of the new constructions in the literature are switching equivalent. It seems that the switching idea is quite powerful to construct new APN functions, since many of the new APN functions listed by Dillon in [9] are within just one switching class. In the case $n=8$, the EA switching class of the Gold function $x^{3}$ contains 17 CCZ inequivalent functions!

\section*{4 Computational results and open problems}
There is, up to equivalence, just one APN mapping $\mathbb{F}_{2}{ }^{n} \rightarrow \mathbb{F}_{2}{ }^{n}$ for $n \leq 4$, hence no interesting things happen in these cases. In the case $n=5$, a complete classification of APN functions (up to CCZ equivalence) is contained in [6]. We summarize our computational results in the following tables. We also include some interesting CCZ invariants:

\begin{itemize}
  \item $\Delta$ - and $\Gamma$-Rank.
  \item Orders of automorphism groups of $\operatorname{dev}\left(D_{F}\right), \operatorname{dev}\left(G_{F}\right)$ and $\mathcal{M}\left(G_{F}\right)$.
  \item Walsh spectrum (see the comments following Proposition 1 about the CCZ invariance of the Walsh spectrum).
\end{itemize}

If all these invariants are the same, we use a direct test to check that the examples are CCZ inequivalent, hence the reader can be sure that all the examples in the following tables are CCZ inequivalent. However, we do not claim that our tables are complete in the sense that they contain all possible CCZ equivalence classes of APN functions with $n=6,7,8$ and 9 .

We list the Walsh spectrum only if it is different from the Walsh spectrum of $x^{3}$. The Walsh spectrum of $x^{3}$ is called classical.

We number the examples as $1.1,1.2, \ldots, 2.1,2.2, \ldots$ etc. The first number describes the EA (CCZ) switching class, and the second number the CCZ inequivalent examples within this class. Our search was complete in the sense that, starting from the known APN functions, we searched through the entire EA (CCZ) switching class. Hence any new APN function must be a member of a new switching class. We used the examples in [22] and [9] as the starting cases.

In the case $n \leq 7$, we searched through the entire CCZ switching class, in the case $n=8,9$ through the EA switching class only.

Each switching class (no matter whether it is the EA switching or the CCZ switching class) may contain several CCZ inequivalent functions. In our tables, we list all CCZ inequivalent examples within one switching class.

Some comments about the sizes of the automorphism groups are in order: The automorphism groups contain the $2^{2 n}$ translations $\tau_{a, b}$, see Theorem 4, therefore we divided the group sizes in our tables by $2^{2 n}$. We were not able to determine the sizes of these groups if $n \geq 8$. However, it was possible to determine the multiplier group $\mathcal{M}\left(G_{F}\right)$ of $\operatorname{dev}\left(G_{F}\right)$ : Using MAGMA, we determined the automorphism groups of the associated extended codes $\mathcal{C}_{F}^{\text {ext }}$. In the cases $n \leq 7$, the automorphism groups of $\operatorname{dev}\left(G_{F}\right)$ have been always the groups generated by the multipliers plus the translations; therefore, it may be possible that the group sizes in Table 10 describe actually the sizes of the full automorphism groups. We did not determine the multiplier groups $\mathcal{M}\left(D_{F}\right)$.

It is quite interesting to look at the automorphism groups of the $\operatorname{designs} \operatorname{dev}\left(G_{F}\right)$ because they give some information about $F$ :

Theorem 10. Let $F(x)$ be an APN function on $\mathbb{F}_{2^{n}}$. Let $v=\left|\mathcal{M}\left(G_{F}\right)\right|$. Then the following holds:\\
(1.) If $F(x)$ is $C C Z$ quadratic, then $2^{n}$ divides $v$.\\
(2.) If $F$ is $C C Z$ equivalent to a power mapping, then $n \cdot\left(2^{n}-1\right)$ divides $v$.\\
(3.) If $F$ is $C C Z$ equivalent to a polynomial in $\mathbb{F}_{2}[x]$, then $n$ divides $v$.

Proof. The first statement is simply Theorem 4. If $F(x) \in \mathbb{F}_{2}[x]$, then the linear mapping $\mathbb{F}_{2^{n}} \times \mathbb{F}_{2^{n}} \rightarrow \mathbb{F}_{2^{n}} \times \mathbb{F}_{2^{n}}$ defined by $(x, y) \mapsto\left(x^{2}, y^{2}\right)$ has order $n$, and it fixes the set $G_{F}$. This shows (3). If $F(x)=\gamma \cdot x^{d}$, then the $2^{n}-1$ mappings defined by $(x, y) \mapsto\left(\alpha x, \alpha^{d} \cdot y\right), \alpha \in \mathbb{F}_{2^{n}}$, fix the set $G_{F}$. Moreover, we may assume that $\gamma=1$, otherwise we replace $F$ by the CCZ equivalent function $\frac{1}{\gamma} F$, which shows (2).

We did mention already that quadratic functions are crooked. The following question is of interest:

Problem 3. Are all crooked functions quadratic?\\
A function defined by $\sum_{i, j} \alpha_{i, j} x^{2^{i}+2^{j}}$ is quadratic. Therefore, Theorems 10 and 7 show the following:

Remark 2. The only non-quadratic functions in Tables 3, 5 and 7 are the function no. 2.1 in Table 3, the function no. 2.12 in Table 5, the functions no. 5.1, 6.1 and 7.1 in Table 7 and the functions no. 4.1, 5.1, 6.1 and 7.1 in Table 11. Moreover, none of these functions is crooked.

Our new function 14.3 in Table 7 is inequivalent to all polynomials with coefficients in $\mathbb{F}_{2}$ (Theorem 10).

We note that Table 3 contains, up to CCZ-equivalence, all APN functions on $\mathbb{F}_{2}{ }^{5}$. In Tables 5 and 7, we determined the CCZ switching classes, starting from the known CCZ inequivalent APN functions. In Tables 9 and 11, we did not apply the switching construction to all the members of the CCZ equivalence classes. In other words, we started with the functions listed in the tables below and determined the EA switching classes, not the CCZ switching classes.

In Table 5, the Example 2.12 is new (see Theorem 11), and in Table 7, the Example 14.3 is new.

In our opinion, the most interesting function is the following non-quadratic example, see also [7]:

Theorem 11. Let $\mathbb{F}_{2^{6}}$ be the finite field which is constructed as the splitting field of $x^{6}+x^{4}+ x^{3}+x+1 \in \mathbb{F}_{2}[x]$. Let $u$ be a root of this polynomial in $\mathbb{F}_{2^{6}}$. Then the function $F: \mathbb{F}_{2^{6}} \rightarrow \mathbb{F}_{2^{6}}$ with


\begin{align*}
F(x)= & x^{3}+u^{17}\left(x^{17}+x^{18}+x^{20}+x^{24}\right)+ \\
& u^{14}\left(\operatorname{tr}\left(u^{52} x^{3}+u^{6} x^{5}+u^{19} x^{7}+u^{28} x^{11}+u^{2} x^{13}\right)+\right.  \tag{9}\\
& \left.\operatorname{tr}_{8 / 2}\left(\left(u^{2} x\right)^{9}\right)+\operatorname{tr}_{4 / 2}\left(x^{21}\right)\right)
\end{align*}


(entry 2.11 in Table 5) is an APN function which is a switching neighbour of Function 2.4 in Table 5. Here $\operatorname{tr}_{8 / 2}$ and $\operatorname{tr}_{4 / 2}$ denote the relative trace $\mathbb{F}_{8} \rightarrow \mathbb{F}_{2}$ and $\mathbb{F}_{4} \rightarrow \mathbb{F}_{2}$. The function cannot be $C C Z$ equivalent to a crooked function. Moreover, it is $C C Z$ inequivalent to a power mapping.

Proof. It is quite easy to check that the function is APN and that it is a switching neighbour of $x^{3}+u^{17}\left(x^{17}+x^{18}+x^{20}+x^{24}\right)$. Note that the $\Delta$-rank $152>2^{7}$ of this function is too big for a crooked function, see Theorem 7. Similarly, the function is inequivalent to a power mapping since the order of the multiplier group is not divisible by $2^{6}-1$, see Theorem 10 .

Two APN functions which are switching equivalent in the narrow sense are quite "similar", they are "almost" equal (see Proposition 3). Therefore, the following question seems natural:

Problem 4. Is there a property which distinguishes switching equivalent functions from those which are not switching equivalent?

\begin{table}[h]
\begin{center}
\begin{tabular}{|l|l|}
\hline
$n$ & $p(x)$ \\
\hline
5 & $x^{5}+x^{2}+1$ \\
\hline
6 & $x^{6}+x^{4}+x^{3}+x+1$ \\
\hline
7 & $x^{7}+x+1$ \\
\hline
8 & $x^{8}+x^{4}+x^{3}+x^{2}+1$ \\
\hline
9 & $x^{9}+x^{4}+1$ \\
\hline
\end{tabular}
\captionsetup{labelformat=empty}
\caption{Table 2: Used primitive polynomials $p(x)$}
\end{center}
\end{table}

Our computational results are quite pessimistic regarding this question. It seems that there is no property of an APN function which is preserved under switching. Also the sizes of the equivalence classes seem to behave "strange": In the $n=8$ case, there is one large switching class, but in the case $n=7$ there are only small classes.

It seems that many inequivalent APN functions exist. So, in our opinion, the main question about APN functions is to determine at least a lower bound for the number of inequivalent ones. Let $\operatorname{APN}(n)$ denote the number of CCZ inequivalent APN functions $\mathbb{F}_{2}{ }^{n} \rightarrow \mathbb{F}_{2}{ }^{n}$. We ask:

Problem 5. Does the function $A P N(n)$ grow exponentially?\\
This paper contains a new non-quadratic APN function. We think that it is worth to search for more examples:

Problem 6. Find more non-quadratic APN functions.\\
We described the switching construction in quite a general form, and then specialized to the case of 1-dimensional subspaces $U \leq\{0\} \times \mathbb{F}_{2}{ }^{n}$. In this case, it was (rather) easy to find switching neighbours. But of course it is also possible to use higher-dimensional subspaces, or even subspaces not contained in $\{0\} \times \mathbb{F}_{2}{ }^{n}$. These cases are more difficult to handle, but we do not think it is impossible. If $U$ becomes larger, the projections are further away from being APN, therefore the "lifting" will become more difficult. On the other hand, if we are using higher-dimensional subspaces $U$, we obtain more freedom for the lifting. But if $U$ is too big, the approach will most likely become useless: If, in the extreme case, $U=\{0\} \times \mathbb{F}_{2}{ }^{n}$, then all APN functions $F$ project onto the same set $\varphi_{U}\left(G_{F}\right)$.

A generalization of the switching idea to functions between fields of odd characteristic is obvious. Therefore, this approach could also be applied to PN functions:

Problem 7. Try to use the switching idea for other subspaces $U$ or for $P N$ functions.\\
Finally, we come back to the original motivation for studying APN functions: APN's are used in cryptography because they are highly nonlinear. Moreover, functions used in cryptography should quite often have large algebraic degree. Therefore, quadratic functions are usually "weak" regarding applications. But our paper shows that functions of large degree can be quite similar to quadratic (i.e. weak) functions via our projection idea. Therefore, it may be worth to see whether "switching" can be used for cryptanalysis.

In the following tables, it is important to know the primitive element that was used to construct the finite fields. In Table 2, we list these polynomials $p(x)$. The primitive element $u$ used later in the tables is a root of $p(x)$ in $\mathbb{F}_{2}[x] /(p(x))$, see [35] for more information about finite field extensions. In the following tables, the notation (No $\mathrm{x} . \mathrm{y})+f(x)$ means that you take the function which is listed in No x.y and add $f(x)$ to it. This notation has the advantage that you see immediately why the functions are switching equivalent. In some cases, we gave a second representative from the respective CCZ class (denoted by $\sim$ ), which is the representative from [9].

\begin{table}[h]
\begin{center}
\begin{tabular}{|l|l|l|}
\hline
\multicolumn{3}{|c|}{$n=5$} \\
\hline\hline
No. & $F(x)$ &  \\
\hline
1.1 & $x^{3}$ &  \\
1.2 & $x^{5}$ & $($ No. 1.1 $)+u^{29} \operatorname{tr}\left(x^{3}+u^{24} x^{5}\right)$ \\
\hline
2.1 & $x^{-1}$ &  \\
\hline
\end{tabular}
\captionsetup{labelformat=empty}
\caption{Table 3: All switching classes of APN's in $\mathbb{F}_{2}^{5}$}
\end{center}
\end{table}

\begin{table}[h]
\begin{center}
\captionsetup{labelformat=empty}
\caption{Table 3: All switching classes of APN's in $\mathbb{F}_{2}^{5}$}
\begin{tabular}{|l|l|l|l|l|l|}
\hline
\multicolumn{6}{|c|}{$n=5$} \\
\hline
No. & $\Gamma$-rank & $\Delta$-rank & $\frac{\left|\operatorname{Aut}\left(\operatorname{dev}\left(G_{F}\right)\right)\right|}{2^{10}}$ & $\frac{\left|\operatorname{Aut}\left(\operatorname{dev}\left(D_{F}\right)\right)\right|}{2^{10}}$ & Walsh spectrum \\
\hline
1.1 & 330 & 42 & $2^{5} \cdot 5 \cdot 31$ & $2^{5} \cdot 5 \cdot 31$ & classical \\
\hline
1.2 & 330 & 42 & $2^{5} \cdot 5 \cdot 31$ & $2^{10} \cdot 5 \cdot 31$ & classical \\
\hline
2.1 & 496 & 232 & $2 \cdot 5 \cdot 31$ & $2 \cdot 5 \cdot 31$ & \( \begin{aligned} & \{* 32[1], 12[21] \\ & 8[195], 4[260] \\ & 0[217],-4[205] \\ & \left.-8[115],-12[10]^{*}\right\} \end{aligned} \) \\
\hline
\end{tabular}
\end{center}
\end{table}

\begin{table}[h]
\begin{center}
\captionsetup{labelformat=empty}
\caption{Table 4: Invariants of switching classes in Table 3}
\begin{tabular}{|l|l|l|}
\hline
\multicolumn{3}{|c|}{$n=6$} \\
\hline
No. & No. in [9] & $F(x)$ \\
\hline
1.1 & 1 & $x^{3}$ \\
\hline
1.2 & 2 & (No. 1.1) $+u\left(\operatorname{tr}\left(u^{56} x^{3}\right)+\operatorname{tr}_{8 / 2}\left(u^{18} x^{9}\right) \quad\left(\sim x^{3}+u^{11} x^{6}+u x^{9}\right)\right.$ \\
\hline
2.1 & 5 & $x^{3}+u x^{24}+x^{10}$ \\
\hline
2.2 & 6 & (No. 2.1) $+u^{3}\left(\operatorname{tr}\left(u^{10} x^{3}+u^{53} x^{5}\right)+\operatorname{tr}_{8 / 2}\left(u^{36} x^{9}\right)\right)$ \\
\hline
2.3 & 9 & $($ No. 2.1$)+\operatorname{tr}\left(u^{34} x^{3}+u^{48} x^{5}\right)+\operatorname{tr}_{8 / 2}\left(u^{9} x^{9}\right)$ \\
\hline
2.4 & 12 & (No. 2.1) $+u^{2}\left(\operatorname{tr}\left(u^{24} x^{3}+u^{28} x^{5}\right)+\operatorname{tr}_{8 / 2}\left(x^{9}\right)\right)$ \\
\hline
2.5 & 3 & $($ No. 2.3$)+u^{42}\left(\operatorname{tr}\left(u^{10} x^{3}+u^{51} x^{5}\right)+\operatorname{tr}_{8 / 2}\left(u^{9} x^{9}\right)\right)$ \\
\hline
2.6 & 7 & (No. 2.3) $+u^{23}\left(\operatorname{tr}\left(u^{31} x^{3}+u^{49} x^{5}\right)+\operatorname{tr}_{8 / 2}\left(u^{9} x^{9}\right)\right)$ \\
\hline
2.7 & 8 & (No. 2.3) $+u^{12}\left(\operatorname{tr}\left(u^{42} x^{3}+u^{13} x^{5}\right)+\operatorname{tr}_{8 / 2}\left(u^{54} x^{9}\right)\right)$ \\
\hline
2.8 & 10 & (No. 2.3) $+u\left(\operatorname{tr}\left(u^{51} x^{3}+u^{60} x^{5}\right)+\operatorname{tr}_{8 / 2}\left(u^{18} x^{9}\right)\right)$ \\
\hline
2.9 & 11 & $($ No. 2.3$)+u^{14}\left(\operatorname{tr}\left(u^{18} x^{3}+u^{61} x^{5}\right)+\operatorname{tr}_{8 / 2}\left(u^{18} x^{9}\right)\right)$ \\
\hline
2.10 & 13 & (No. 2.3) $+u^{17}\left(\operatorname{tr}\left(u^{50} x^{3}+u^{56} x^{5}\right)\right)$ \\
\hline
2.11 & new (Theorem 11) & $($ No. 2.3 $)+u^{19}\left(\operatorname{tr}\left(u^{11} x^{3}+u^{7} x^{5}+u^{38} x^{7}+u^{61} x^{11}+u^{23} x^{13}\right)+\right. \left.\operatorname{tr}_{8 / 2}\left(u^{54} x^{9}\right)+\operatorname{tr}_{4 / 2}\left(u^{42} x^{21}\right)\right)$ \\
\hline
2.12 & 4 & (No. 2.4) $+u\left(\operatorname{tr}\left(u^{54} x^{3}+u^{47} x^{5}\right)+\operatorname{tr}_{8 / 2}\left(u^{9} x^{9}\right)\right)$ \\
\hline
\end{tabular}
\end{center}
\end{table}

Table 5: Known switching classes of APN's in $\mathbb{F}_{2}^{6}$

\begin{table}[h]
\begin{center}
\begin{tabular}{|l|l|l|l|l|l|}
\hline
\multicolumn{6}{|c|}{$n=6$} \\
\hline
No. & $\Gamma$-rank & $\Delta$-rank & $\frac{\left|\operatorname{Aut}\left(\operatorname{dev}\left(G_{F}\right)\right)\right|}{2^{12}}$ & $\frac{\left|\operatorname{Aut}\left(\operatorname{dev}\left(D_{F}\right)\right)\right|}{2^{12}}$ & Walsh spectrum \\
\hline
1.1 & 1102 & 94 & $2^{7} \cdot 3^{3} \cdot 7$ & $2^{8} \cdot 3^{3} \cdot 7$ & classical \\
\hline
1.2 & 1146 & 94 & $2^{6} \cdot 3^{2} \cdot 7$ & $2^{7} \cdot 3^{2} \cdot 7$ & classical \\
\hline
2.1 & 1166 & 96 & $2^{7} \cdot 7$ & $2^{7} \cdot 7$ & classical \\
\hline
2.2 & 1168 & 96 & $2^{6}$ & $2^{6}$ & classical \\
\hline
2.3 & 1170 & 96 & $2^{6}$ & $2^{6}$ & classical \\
\hline
2.4 & 1172 & 96 & $2^{6}$ & $2^{6}$ & classical \\
\hline
2.5 & 1158 & 96 & $2^{6} \cdot 5$ & $2^{6} \cdot 5$ & classical \\
\hline
2.6 & 1170 & 96 & $2^{6} \cdot 5$ & $2^{6} \cdot 5$ & \( \begin{gathered} \{* 64[1], 32[3], \\ 16[160], 8[1656], \\ 0[891],-8[1288], \\ -16[96],-32[1] \quad *\} \end{gathered} \) \\
\hline
2.7 & 1170 & 96 & $2^{6}$ & $2^{6}$ & classical \\
\hline
2.8 & 1170 & 96 & $2^{6}$ & $2^{6}$ & classical \\
\hline
2.9 & 1172 & 96 & $2^{6}$ & $2^{6}$ & classical \\
\hline
2.10 & 1174 & 96 & $2^{6}$ & $2^{6}$ & classical \\
\hline
2.11 & 1300 & 152 & $2^{3}$ & $2^{3}$ & classical \\
\hline
2.12 & 1166 & 94 & $2^{6} \cdot 7$ & $2^{7} \cdot 7$ & classical \\
\hline
\end{tabular}
\captionsetup{labelformat=empty}
\caption{Table 6: Known switching classes of APN's in $\mathbb{F}_{2}{ }^{6}$ : Invariants}
\end{center}
\end{table}

\begin{table}[h]
\begin{center}
\begin{tabular}{|l|l|l|}
\hline
\multicolumn{3}{|c|}{$n=7$} \\
\hline
No. & No. in [9] & $F(x)$ \\
\hline
1.1 & 1 & $x^{3}$ \\
\hline
1.2 & 7 & $x^{3}+\operatorname{tr}\left(x^{9}\right)$ \\
\hline
2.1 & 8 & \( \begin{aligned} & x^{34}+x^{18}+x^{5} \\ & (\mathrm{No.} 2.1)+\operatorname{tr}\left(x^{3}+u^{103} x^{5}+u^{5} x^{9}\right) \\ & \left(\sim x^{3}+x^{17}+x^{33}+x^{34}\right) \end{aligned} \) \\
\hline
3.1 & 3 & $x^{5}$ \\
\hline
4.1 & 2 & $x^{9}$ \\
\hline
5.1 & 4 & $x^{13}$ \\
\hline
6.1 & 5 & $x^{57}$ \\
\hline
7.1 & 6 & $x^{-1}$ \\
\hline
8.1 & 9 & $x^{65}+x^{10}+x^{3}$ \\
\hline
9.1 & 13 & $x^{3}+x^{9}+x^{18}+x^{66}$ \\
\hline
 &  & \( \begin{aligned} & x^{3}+x^{12}+x^{17}+x^{33} \\ & (\text { No. } 10.1)+\operatorname{tr}\left(u^{41} x^{3}+u^{17} x^{5}+u^{34} x^{9}\right) \\ & \left(\sim x^{3}+x^{17}+x^{20}+x^{34}+x^{66}\right) \end{aligned} \) \\
\hline
11.1 & 15 & $x^{3}+x^{20}+x^{34}+x^{66}$ \\
\hline
12.1 & 16 & $x^{3}+x^{12}+x^{40}+x^{72}$ \\
\hline
13.1 & 12 & $x^{3}+x^{5}+x^{10}+x^{33}+x^{34}$ \\
\hline
14.2 & 17 & \( \begin{aligned} & x^{3}+x^{6}+x^{34}+x^{40}+x^{72} \\ & (\text { No. } 14.1)+\operatorname{tr}\left(u^{105} x^{3}+u^{84} x^{5}+u^{123} x^{9}\right) \\ & \left(\sim x^{3}+x^{5}+x^{6}+x^{12}+x^{33}+x^{34}\right) \\ & (\text { No. } 14.1)+u^{27}\left(\operatorname{tr}\left(u^{20} x^{3}+u^{94} x^{5}+u^{66} x^{9}\right)\right) \end{aligned} \) \\
\hline
\end{tabular}
\captionsetup{labelformat=empty}
\caption{Table 7: Known switching classes of APN's in $\mathbb{F}_{2}^{7}$}
\end{center}
\end{table}

\begin{table}[h]
\begin{center}
\begin{tabular}{|l|l|l|l|l|l|}
\hline
\multicolumn{6}{|c|}{$n=7$} \\
\hline
No. & $\Gamma$-rank & $\Delta$-rank & $\frac{\left|\operatorname{Aut}\left(\operatorname{dev}\left(G_{F}\right)\right)\right|}{2^{14}}$ & $\frac{\left|\operatorname{Aut}\left(\operatorname{dev}\left(D_{F}\right)\right)\right|}{2^{14}}$ & Walsh spectrum \\
\hline
1.1 & 3610 & 198 & $2^{7} \cdot 7 \cdot 127$ & $2^{7} \cdot 7 \cdot 127$ & classical \\
\hline
1.2 & 4026 & 212 & $2^{7} \cdot 7$ & $2^{7} \cdot 7$ & classical \\
\hline
2.1 & 4034 & 210 & $2^{7} \cdot 7$ & $2^{7} \cdot 7$ & classical \\
\hline
2.2 & 4040 & 212 & $2^{7} \cdot 7$ & $2^{7} \cdot 7$ & classical \\
\hline
3.1 & 3708 & 198 & $2^{7} \cdot 7 \cdot 127$ & $2^{7} \cdot 7 \cdot 127$ & classical \\
\hline
4.1 & 3610 & 198 & $2^{7} \cdot 7 \cdot 127$ & $2^{14} \cdot 7 \cdot 127$ & classical \\
\hline
5.1 & 4270 & 338 & $7 \cdot 127$ & $7 \cdot 127$ & classical \\
\hline
6.1 & 4704 & 436 & $7 \cdot 127$ & $7 \cdot 127$ & classical \\
\hline
7.1 & 8128 & 4928 & $2 \cdot 7 \cdot 127$ & $2 \cdot 7 \cdot 127$ & \{* 128[1], 20[1029], 16[1498], 12[1526], 8[1897], 4[1827], 0[2032], $-4[1729],-8[1659]$, -12[1268], -16[1169], -20[749] *\} \\
\hline
8.1 & 4038 & 212 & $2^{7} \cdot 7$ & $2^{7} \cdot 7$ & classical \\
\hline
9.1 & 4044 & 212 & $2^{7} \cdot 7$ & $2^{7} \cdot 7$ & classical \\
\hline
10.1 & 4048 & 210 & $2^{7} \cdot 7$ & $2^{7} \cdot 7$ & classical \\
\hline
10.2 & 4040 & 210 & $2^{7} \cdot 7$ & $2^{7} \cdot 7$ & classical \\
\hline
11.1 & 4048 & 210 & $2^{7} \cdot 7$ & $2^{7} \cdot 7$ & classical \\
\hline
12.1 & 4048 & 210 & $2^{7} \cdot 7$ & $2^{7} \cdot 7$ & classical \\
\hline
13.1 & 4040 & 212 & $2^{7} \cdot 7$ & $2^{7} \cdot 7$ & classical \\
\hline
14.1 & 4048 & 212 & $2^{7} \cdot 7$ & $2^{7} \cdot 7$ & classical \\
\hline
14.2 & 4050 & 210 & $2^{7} \cdot 7$ & $2^{7} \cdot 7$ & classical \\
\hline
14.3 & 4046 & 212 & $2^{7}$ & $2^{7}$ & classical \\
\hline
\end{tabular}
\captionsetup{labelformat=empty}
\caption{Table 8: Known switching classes of APN's in $\mathbb{F}_{2}{ }^{7}$ : Invariants}
\end{center}
\end{table}

\begin{table}[h]
\begin{center}
\begin{tabular}{|l|l|l|}
\hline
\multicolumn{3}{|c|}{$n=8$} \\
\hline
No. & No. in [9] & F(x) \\
\hline
1.1 & 1 & $x^{3}$ \\
\hline
1.2 & 5 & (No. 1.1) $+\operatorname{tr}\left(u^{48} x^{3}+x^{9}\right) \quad\left(\sim x^{3}+\operatorname{tr} x^{9}\right)$ \\
\hline
1.3 & 7 & (No. 1.1) $+u\left(\operatorname{tr}\left(u^{63} x^{3}+u^{252} x^{9}\right)\right)$ \\
\hline
 &  & $\left(\sim x^{3}+u^{245} x^{33}+u^{183} x^{66}+u^{21} x^{144}\right)$ \\
\hline
1.4 & new & (No. 1.2) $+u^{38}\left(\operatorname{tr}\left(u^{84} x^{3}+u^{213} x^{9}\right)\right)$ \\
\hline
1.5 & new & (No. 1.2) $+u^{51}\left(\operatorname{tr}\left(u^{253} x^{3}+u^{102} x^{9}\right)\right)$ \\
\hline
1.6 & new & (No. 1.3) $+u^{154}\left(\operatorname{tr}\left(u^{68} x^{3}+u^{235} x^{9}\right)\right)$ \\
\hline
1.7 & new & (No. 1.4) $+u^{69}\left(\operatorname{tr}\left(u^{147} x^{3}+u^{20} x^{9}\right)\right)$ \\
\hline
1.8 & new & (No. 1.5) $+u^{68}\left(\operatorname{tr}\left(u^{153} x^{3}+u^{51} x^{9}\right)\right)$ \\
\hline
1.9 & new & (No. 1.6) $+u^{35}\left(\operatorname{tr}\left(u^{216} x^{3}+u^{116} x^{9}\right)\right)$ \\
\hline
1.10 & new & (No. 1.7) $+u^{22}\left(\operatorname{tr}\left(u^{232} x^{3}+u^{195} x^{9}\right)\right)$ \\
\hline
1.11 & 6 & $($ No. 1.8 $)+u^{85}\left(\operatorname{tr}\left(u^{243} x^{3}+u^{170} x^{9}\right)\right) \quad\left(\sim x^{9}+\operatorname{tr}\left(x^{3}\right)\right)$ \\
\hline
\multirow{2}{*}{1.12} & 8 & (No. 1.9) $+u^{103}\left(\operatorname{tr}\left(u^{172} x^{3}+u^{31} x^{9}\right)\right)$ \\
\hline
 &  & $\left(\sim x^{3}+u^{65} x^{18}+u^{120} x^{66}+u^{135} x^{144}\right)$ \\
\hline
1.13 & new & (No. 1.10) $+u^{90}\left(\operatorname{tr}\left(u^{87} x^{3}+u^{141} x^{5}+u^{20} x^{9}\right)+\right. \left.\operatorname{tr}_{16 / 2}\left(u^{51} x^{17}\right)\right)$ \\
\hline
1.14 & new & (No. 1.11) $+u^{5}\left(\operatorname{tr}\left(u^{160} x^{3}+u^{250} x^{9}\right)\right)$ \\
\hline
1.15 & 2 & $($ No. 1.11 $)+u^{102}\left(\operatorname{tr}\left(u^{6} x^{3}+u^{119} x^{9}\right)\right) \quad\left(\sim x^{9}\right)$ \\
\hline
1.16 & new & (No. 1.14) $+u^{64}\left(\operatorname{tr}\left(u^{133} x^{3}+u^{30} x^{9}\right)\right)$ \\
\hline
1.17 & new & (No. 1.16) $+u^{78}\left(\operatorname{tr}\left(u^{235} x^{3}+u^{146} x^{9}\right)\right)$ \\
\hline
2.1 & 4 & $x^{3}+x^{17}+u^{16}\left(x^{18}+x^{33}\right)+u^{15} x^{48}$ \\
\hline
3.1 & 9 & $x^{3}+u^{24} x^{6}+u^{182} x^{132}+u^{67} x^{192}$ \\
\hline
4.1 & 10 & $x^{3}+x^{6}+x^{68}+x^{80}+x^{132}+x^{160}$ \\
\hline
5.1 & 11 & $x^{3}+x^{5}+x^{18}+x^{40}+x^{66}$ \\
\hline
6.1 & 12 & $x^{3}+x^{12}+x^{40}+x^{66}+x^{130}$ \\
\hline
7.1 & 3 & $x^{57}$ \\
\hline
\end{tabular}
\captionsetup{labelformat=empty}
\caption{Table 9: Known switching classes of APN's in $\mathbb{F}_{2}^{8}$}
\end{center}
\end{table}

\begin{table}[h]
\begin{center}
\begin{tabular}{|l|l|l|l|l|}
\hline
\multicolumn{5}{|c|}{$n=8$} \\
\hline
No. & $\Gamma$-rank & $\Delta$-rank & $\left.\mid \mathcal{M}\left(G_{F}\right)\right) \mid$ & Walsh spectrum \\
\hline
1.1 & 11818 & 420 & $2^{11} \cdot 255$ & classical \\
\hline
1.2 & 13800 & 432 & $2^{11} \cdot 3$ & classical \\
\hline
1.3 & 13842 & 436 & $2^{10} \cdot 3$ & classical \\
\hline
1.4 & 14034 & 438 & $2^{8} \cdot 3$ & classical \\
\hline
1.5 & 14032 & 438 & $2^{10} \cdot 3$ & classical \\
\hline
1.6 & 14036 & 438 & $2^{10} \cdot 3$ & classical \\
\hline
1.7 & 14036 & 438 & $2^{9} \cdot 3$ & classical \\
\hline
1.8 & 14032 & 438 & $2^{10} \cdot 3$ & classical \\
\hline
1.9 & 14034 & 438 & $2^{10} \cdot 3$ & classical \\
\hline
1.10 & 14030 & 438 & $2^{9} \cdot 3$ & classical \\
\hline
1.11 & 13804 & 434 & $2^{11} \cdot 3$ & classical \\
\hline
1.12 & 13848 & 438 & $2^{10} \cdot 3$ & classical \\
\hline
1.13 & 14046 & 454 & $2^{9}$ & classical \\
\hline
1.14 & 14036 & 438 & $2^{8} \cdot 3$ & classical \\
\hline
1.15 & 12370 & 420 & $2^{11} \cdot 255$ & classical \\
\hline
1.16 & 14032 & 438 & $2^{9} \cdot 3$ & classical \\
\hline
1.17 & 14028 & 438 & $2^{9} \cdot 3$ & classical \\
\hline
2.1 & 13200 & 414 & $2^{10} \cdot 3^{2} \cdot 5$ & classical \\
\hline
3.1 & 14024 & 438 & $2^{10} \cdot 3$ & classical \\
\hline
4.1 & 14040 & 454 & $2^{11}$ & classical \\
\hline
5.1 & 14044 & 446 & $2^{11}$ & classical \\
\hline
6.1 & 14046 & 438 & $2^{11}$ & classical \\
\hline
7.1 & 15358 & 960 & $2^{3} \cdot 255$ & classical \\
\hline
\end{tabular}
\captionsetup{labelformat=empty}
\caption{Table 10: Known switching classes of APN's in $\mathbb{F}_{2}^{8}$ : Invariants}
\end{center}
\end{table}

\begin{table}[h]
\begin{center}
\begin{tabular}{|l|l|}
\hline
\multicolumn{2}{|c|}{$n=9$} \\
\hline
No. & F(x) \\
\hline
1.1 & $x^{3}$ \\
\hline
1.2 & (No. 1.1) $+\operatorname{tr}\left(x^{9}\right)$ \\
\hline
1.3 & (No. 1.2) $+u^{73} \operatorname{tr}\left(u^{426} x^{3}+u^{292} x^{9}\right)$ \\
\hline
1.4 & (No. 1.2) $+u^{219} \operatorname{tr}\left(u^{303} x^{3}+u^{365} x^{9}\right)$ \\
\hline
2.1 & $x^{5}$ \\
\hline
3.1 & $x^{17}$ \\
\hline
4.1 & $x^{13}$ \\
\hline
5.1 & $x^{19}$ \\
\hline
6.1 & $x^{241}$ \\
\hline
7.1 & $x^{-1}$ \\
\hline
8.1 & $x^{3}+x^{10}+u^{438} x^{136}$ \\
\hline
\end{tabular}
\captionsetup{labelformat=empty}
\caption{Table 11: Known switching classes of APN's in $\mathbb{F}_{2}{ }^{9}$}
\end{center}
\end{table}

\begin{table}[h]
\begin{center}
\begin{tabular}{|l|l|l|l|l|}
\hline
\multicolumn{5}{|c|}{$n=9$} \\
\hline
No. & $\Gamma$-rank & $\Delta$-rank & $\left.\mid \mathcal{M}\left(G_{F}\right)\right) \mid$ & Walsh spectrum \\
\hline
1.1 & 38470 & 872 & $9 \cdot 2^{9} \cdot 511$ & classical \\
\hline
1.2 & 47890 & 920 & $9 \cdot 2^{9}$ & classical \\
\hline
1.3 & 48428 & 930 & $9 \cdot 2^{9}$ & classical \\
\hline
1.4 & 48460 & 944 & $9 \cdot 2^{9}$ & classical \\
\hline
2.1 & 41494 & 872 & $9 \cdot 2^{9} \cdot 511$ & classical \\
\hline
3.1 & 38470 & 872 & $9 \cdot 2^{9} \cdot 511$ & classical \\
\hline
4.1 & 58676 & 3086 & $9 \cdot 511$ & classical \\
\hline
5.1 & 60894 & 3956 & $9 \cdot 511$ & classical \\
\hline
6.1 & 61726 & 3482 & 9.511 & classical \\
\hline
7.1 & 130816 & 93024 & $2 \cdot 9 \cdot 511$ & \{ ${ }^{*} 512$ [1], 44[3327], 40[7434], 36[7389], 32[12231], 28[16965], 24[8019], 20[11952], 16[18963], 12[11772], 8[16353], 4[16467], 0[10220], -4[16237], $-8[15840], \quad-12[11223], \quad-16[17829], \quad-20[11043]$, $-24[7311], \quad-28[15228], \quad-32[10764], \quad-36[6408]$, -40[6363], -44[2805] *\} \\
\hline
8.1 & 48608 & 938 & $3 \cdot 7 \cdot 2^{9}$ & classical \\
\hline
\end{tabular}
\captionsetup{labelformat=empty}
\caption{Table 12: Known switching classes of APN's in $\mathbb{F}_{2}^{9}$ : Invariants}
\end{center}
\end{table}

\section*{References}
[1] T. D. Bending and D. Fon-Der-Flaass, Crooked functions, bent functions, and distance regular graphs, Electron. J. Combin., 5 (1998), pp. Research Paper 34, 14 pp. (electronic).\\[0pt]
[2] T. Beth, D. Jungnickel and H. Lenz, Design Theory, Cambridge University Press, Cambridge, 2 ed., 1999.\\[0pt]
[3] J. Bierbrauer and G. M. Kyureghyan, Crooked binomials, Des. Codes Cryptogr., 46 (2008), pp. 269-301.\\[0pt]
[4] W. Bosma, J. Cannon and C. Playoust, The magma algebra system. i. the user language, J. Symbolic Comput., 24 (1997), pp. 235-265.\\[0pt]
[5] C. Bracken, E. Byrne, N. Markin and G. McGuire, Quadratic almost perfect nonlinear functions with many terms. IACR Cryptology ePrint Archive: 2007/115, 2007.\\[0pt]
[6] M. Brinkmann and G. Leander, On the classification of APN functions up to dimension five, in Abstract Book of the Workshop on coding and cryptography, N. S. D. Augo and J.-P. Tillich, eds., INRIA, 2007, pp. 39-48.\\
$[7]$-, On the classification of APN functions up to dimension five, Des., Codes, Cryptogr., (2008).\\[0pt]
[8] A. E. Brouwer and L. M. G. M. Tolhuizen, A sharpening of the Johnson bound for binary linear codes and the nonexistence of linear codes with Preparata parameters, Des. Codes Cryptogr., 3 (1993), pp. 95-98.\\[0pt]
[9] K. Browning, J. Dillon, R. Kibler and M. McQuistan, APN polynomials and related codes. submitted, 2008.\\[0pt]
[10] L. Budaghyan and C. Carlet, Classes of quadratic APN trinomials and hexanomials and related structures, IEEE \href{http://Trans.Inf.Th}{Trans.Inf.Th}., 54 (2008), pp. 2354-2357.\\[0pt]
[11] Constructing new APN functions from known ones. to appear in Finite Fields Appl., 2009.\\[0pt]
[12] -, On inequivalence between known power APN functions, in International Conference on Boolean Functions: Cryptography and Applications, 2008.\\[0pt]
[13] L. Budaghyan, C. Carlet and G. Leander, Two classes of quadratic APN binomials inequivalent to power functions, IEEE \href{http://Trans.Inf.Th}{Trans.Inf.Th}., 54 (2008), pp. 4218-4229.\\[0pt]
[14] L. Budaghyan, C. Carlet and A. Pott, New classes of almost bent and almost perfect nonlinear polynomials, IEEE Trans. Inform. Theory, 52 (2006), pp. 1141-1152.\\[0pt]
[15] L. Budaghyan and T. Helleseth, New perfect nonlinear monomials over $\mathbb{F}_{p^{2 k}}$ for any odd prime $p$. to be presented at SETA '08.\\[0pt]
[16] E. Byrne, C. Bracken, N. Markin and G. McGuire, New families of quadratic almost perfect nonlinear trinomials and multinomials. preprint, available online at: \href{http://mathsci.ucd.ie/~gmg/}{http://mathsci.ucd.ie/\~{}gmg/}, 2007.\\[0pt]
[17] E. Byrne and G. McGuire, Certain new quadratic APN functions are not APN infinitely often, in Abstract Book of the Workshop on coding and cryptography, N. S. D. Augo and J.-P. Tillich, eds., INRIA, 2007, pp. 59-68.\\[0pt]
[18] A. Canteaut, P. Charpin and H. Dobbertin, Weight divisibility of cyclic codes, highly nonlinear functions on $\mathbf{F}_{2^{m}}$, and crosscorrelation of maximum-length sequences, SIAM J. Discrete Math., 13 (2000), pp. 105-138.\\[0pt]
[19] C. Carlet, P. Charpin and V. Zinoviev, Codes, bent functions and permutations suitable for DES-like cryptosystems, Des. Codes Cryptogr., 15 (1998), pp. 125-156.\\[0pt]
[20] F. Chabaud and S. Vaudenay, Links between differential and linear cryptanalysis, in Advances in Cryptology - EUROCRYPT 94, A. D. Santis, ed., vol. 950 of Lecture Notes in Computer Science, New York, 1995, Springer-Verlag, pp. 356-365.\\[0pt]
[21] J. Dillon and H. Dobbertin, New cyclic difference sets with Singer parameters., Finite Fields Appl., 10 (2004), pp. 342-389.\\[0pt]
[22] J. F. Dillon. slides from talk given at "Polynomials over Finite Fields and Appliocations", held at Banff International Research Station, 2006.\\[0pt]
[23] C. Ding and J. Yuan, A new family of skew Paley-Hadamard difference sets, J. Comb. Theory Ser.A, 113 (2006), pp. 1526-1535.\\[0pt]
[24] H. Dobbertin, Almost perfect nonlinear power functions on $G F\left(2^{n}\right)$ : The Niho case, Information and Computation, 151 (1999), pp. 57-72.\\[0pt]
[25] -, Almost perfect nonlinear power functions on $\mathrm{GF}\left(2^{n}\right)$ : the Welch case, IEEE Trans. Inform. Theory, 45 (1999), pp. 1271-1275.\\[0pt]
[26] -, Almost perfect nonlinear power functions on $G F\left(2^{n}\right)$ : A new case for $n$ divisible by 5 , in Proceedings of the conference on Finite Fields and Applications, Augsburg 1999, D. Jungnickel and H. Niederreiter, eds., Berlin, 2001, Springer-Verlag, pp. 113-121.\\[0pt]
[27] Y. Edel, G. Kyureghyan and A. Pott, A new APN function which is not equivalent to a power mapping, IEEE Trans. Inform. Theory, 52 (2006), pp. 744-747.\\[0pt]
[28] R. Gold, Maximal recursive sequences with 3 -valued recursive cross-correlation function, IEEE Trans. Inf. Th., 14 (1968), pp. 154-156.\\[0pt]
[29] S. W. Golomb and G. Gong, Signal design for good correlation, Cambridge University Press, Cambridge, 2005. For wireless communication, cryptography, and radar.\\[0pt]
[30] F. Göloğlu and A. Pott, Almost perfect nonlinear functions: A possible geometric approach (25 pages). Presented at Academic Contact Forum Coding Theory and Cryptography, Brussels, 2008.\\[0pt]
[31] H. Janwa and R. M. Wilson, Hyperplane sections of Fermat varieties in $\mathbf{P}^{3}$ in char. 2 and some applications to cyclic codes, in Applied algebra, algebraic algorithms and error-correcting codes (San Juan, PR, 1993), vol. 673 of Lecture Notes in Comput. Sci., Springer, Berlin, 1993, pp. 180-194.\\[0pt]
[32] T. Kasami, The weight enumerators for several classes of subcodes of the 2 nd order binary Reed-Muller codes, Information and Control, 18 (1971), pp. 369-394.\\[0pt]
[33] G. M. Kyureghyan, Crooked maps in $\mathbb{F}_{2^{n}}$, Finite Fields Appl., 13 (2007), pp. 713-726.\\[0pt]
[34] G. Lachaud and J. Wolfmann, The weights of the orthogonals of the extended quadratic binary Goppa codes, IEEE Trans. Inform. Theory, 36 (1990), pp. 686-692.\\[0pt]
[35] R. Lidl and H. Niederreiter, Finite Fields, vol. 20 of Encyclopedia of Mathematics and its Applications, Cambridge University Press, 2nd ed., 1997.\\[0pt]
[36] K. Nyberg, Differentially uniform mappings for cryptography, in Advances in Cryptography. EUROCRYPT'93, vol. 765 of Lecture Notes in Computer Science, New York, 1994, SpringerVerlag, pp. 55-64.\\[0pt]
[37] Z. Zha, G. M. Kyureghyan and X. Wang, A new family of perfect nonlinear binomials. submitted, 2008.


\end{document}